%! Characteristics of Logarithmic Functions — Unit 8, Lesson 8.0 — Group Activity (Answer Key)
\documentclass[10pt]{article}
\usepackage{algebra2-article}
\usepackage{algebra2-key}   % pulls in -boxes; do NOT also load -boxes
\newcommand{\TallMath}[1]{$\displaystyle #1\rule[-1.4em]{0pt}{3.2em}$}
\newcommand{\lgrid}[4]{%
  \draw[step=1, linegray!40, very thin] (#1,#3) grid (#2,#4);
  \draw[->, semithick, charcoal] (#1,0)--(#2,0) node[below right, font=\scriptsize]{$x$};
  \draw[->, semithick, charcoal] (0,#3)--(0,#4) node[left, font=\scriptsize]{$y$};}
% pgfmath has no log_b, so every logarithmic curve is drawn by parameterizing on y:
%   y = log_b(x-h) + k   <=>   x = b^(y-k) + h.   Args: {b}{h}{k}{ymin}{ymax}
\newcommand{\logcurve}[5]{\draw[very thick, forest, domain=#4:#5, samples=120, smooth]
  plot ({pow(#1,\x-(#3))+(#2)},{\x});}

\begin{document}

\pageheader{Unit 8, Lesson 8.0}{Group Activity --- Answer Key}
\namepartnerperiod

\vspace{0.08in}

\begin{headlinebox}{forestbg}
\small\textbf{Reading Logarithmic Graphs.} A logarithmic function tells you \emph{where it is allowed
to exist} before it tells you anything else --- and the \textbf{argument} is what decides. With your
partner, read each graph carefully and \emph{justify} your answers. Write every domain and range in
\textbf{interval notation}, and check each one against the argument. Watch your brackets: a
logarithm's boundary is an \textbf{asymptote}, not an endpoint, so it is \emph{never} included.
\end{headlinebox}

\vspace{0.06in}

\begin{tcolorbox}[colback=white, colframe=black!40, title=\textbf{Tier R --- Remediate},
    fonttitle=\bfseries, arc=1.5mm, left=3mm, right=3mm, top=2mm, bottom=2mm, breakable]
The graph shows $g(x) = \log_3 x$.
\begin{center}
\begin{tikzpicture}[scale=0.46]
  \lgrid{-3}{11}{-3}{3}
  \draw[navy, dashed, semithick] (0,-3)--(0,3);
  \begin{scope}
    \clip (-3,-3) rectangle (11,3);
    \logcurve{3}{0}{0}{-3}{2.19}
  \end{scope}
  \fill[charcoal] (1,0) circle (2.6pt) node[below right, font=\scriptsize]{$(1,0)$};
  \fill[charcoal] (3,1) circle (2.6pt) node[above left, font=\scriptsize]{$(3,1)$};
  \fill[charcoal] (9,2) circle (2.6pt) node[above left, font=\scriptsize]{$(9,2)$};
\end{tikzpicture}
\end{center}
\begin{enumerate}[leftmargin=1.7em, itemsep=7pt]
  \item The \textbf{argument} is \ans{$x$}. Set it $>0$: \ans{$x>0$}, so the
        \textbf{domain} is \ans{$(0,\infty)$}.
  \item The graph leans on a vertical line it never touches. That line is the \textbf{vertical
        asymptote}: \ans{$x=0$}. \\
        Why can the graph never reach it? \ans{because $x=0$ is not in the domain --- no power of 3 equals 0}
  \item \textbf{Range:} \ans{$(-\infty,\infty)$} \quad \textbf{$x$-intercept:} \ans{$(1,0)$}
  \item Try to find the \textbf{$y$-intercept} by substituting $x=0$. You would need to answer
        ``$3$ to what power gives \ans{0}?'', which is \ans{impossible}. \\
        What does that tell you about whether this graph has a $y$-intercept? \ans{it has none --- and never will}
  \item $g$ is \textbf{increasing} on \ans{$(0,\infty)$}. \quad Its absolute extrema:
        \ans{none}, because the graph falls forever on one side and climbs forever on the other.
  \item \textbf{Positive / negative:} $g(x)<0$ on \ans{$(0,1)$}; \quad $g(x)>0$ on
        \ans{$(1,\infty)$}.
  \item \textbf{End behavior:} as $x\to+\infty$, $g(x)\to$ \ans{$+\infty$}. \\
        On the left the input cannot run to $-\infty$ --- careful! Fill in the correct notation:
        as $x\to$ \ans{$0^{+}$}, $g(x)\to$ \ans{$-\infty$}.
\end{enumerate}
\end{tcolorbox}

\begin{tcolorbox}[colback=white, colframe=black!40, title=\textbf{Tier A --- Approaching Proficiency},
    fonttitle=\bfseries, arc=1.5mm, left=3mm, right=3mm, top=2mm, bottom=2mm, breakable]
Two logarithmic functions --- and \emph{only one of them has a $y$-intercept}. Fill in the whole
table, then answer the two justification questions.
\begin{center}
\begin{tikzpicture}[scale=0.42]
  % p(x) = log_2(x+2)
  \begin{scope}
    \lgrid{-4}{8}{-4}{4}
    \draw[navy, dashed, semithick] (-2,-4)--(-2,4);
    \begin{scope}
      \clip (-4,-4) rectangle (8,4);
      \logcurve{2}{-2}{0}{-4}{3.33}
    \end{scope}
    \fill[charcoal] (-1,0) circle (3pt);
    \fill[charcoal] (0,1) circle (3pt);
    \fill[charcoal] (2,2) circle (3pt);
    \fill[charcoal] (6,3) circle (3pt);
    \node[charcoal, font=\scriptsize] at (2,-5.4) {$p(x)=\log_2(x+2)$};
  \end{scope}
  % q(x) = log_{1/2}(x)
  \begin{scope}[xshift=16cm]
    \lgrid{-2}{10}{-4}{4}
    \draw[navy, dashed, semithick] (0,-4)--(0,4);
    \begin{scope}
      \clip (-2,-4) rectangle (10,4);
      \logcurve{0.5}{0}{0}{-3.33}{4}
    \end{scope}
    \fill[charcoal] (1,0) circle (3pt);
    \fill[charcoal] (2,-1) circle (3pt);
    \fill[charcoal] (4,-2) circle (3pt);
    \fill[charcoal] (8,-3) circle (3pt);
    \node[charcoal, font=\scriptsize] at (4,-5.4) {$q(x)=\log_{1/2} x$};
  \end{scope}
\end{tikzpicture}
\end{center}
\vspace{0.20in}
\renewcommand{\arraystretch}{1.5}
\begin{tabularx}{\linewidth}{@{}>{\bfseries}p{0.30\linewidth} X X@{}}
 & \textbf{$p(x)=\log_2(x+2)$} & \textbf{$q(x)=\log_{1/2} x$} \\
\hline
Base: $b>1$ or $0<b<1$? & \ans{$b>1$} & \ans{$0<b<1$} \\
Argument & \ans{$x+2$} & \ans{$x$} \\
Domain & \ans{$(-2,\infty)$} & \ans{$(0,\infty)$} \\
Range & \ans{all reals} & \ans{all reals} \\
Vertical asymptote & \ans{$x=-2$} & \ans{$x=0$} \\
$x$-intercept & \ans{$(-1,0)$} & \ans{$(1,0)$} \\
$y$-intercept & \ans{$(0,1)$} & \ans{none} \\
Increasing or decreasing? & \ans{increasing} & \ans{decreasing} \\
As $x\to+\infty$ & \ans{$y\to+\infty$} & \ans{$y\to-\infty$} \\
\end{tabularx}
\vspace{4pt}

\textbf{Justify 1.} $q$ has no $y$-intercept, but $p$ does. Point to the exact feature of each
\emph{equation} that causes this --- and explain why that feature has the effect it does.
\ansline{A $y$-intercept requires $x=0$ to be a legal input. In $q$ the argument is just $x$, so the
domain is $x>0$ and $x=0$ sits exactly on the asymptote --- excluded, so there is no $y$-intercept.\\[3pt]
In $p$ the argument is $x+2$, so the domain is $x>-2$, and $0$ is comfortably inside it:
$p(0)=\log_2 2=1$, giving the $y$-intercept $(0,1)$. The ``$+2$'' pushed the asymptote off the
$y$-axis, and that is what made room for the intercept.}

\textbf{Justify 2.} Both graphs have a vertical asymptote, but they are different lines. Explain how
you could find each one from the \emph{equation alone}, without looking at either picture.
\ansline{Set the \emph{argument} equal to $0$ and solve --- that input is the one the domain stops
at, so it is the asymptote. For $p$: $x+2=0 \Rightarrow x=-2$. For $q$: $x=0$. The domain is then
everything on the side of that line where the argument is \emph{positive}: $(-2,\infty)$ and
$(0,\infty)$.}
\end{tcolorbox}

\begin{tcolorbox}[colback=white, colframe=black!40, title=\textbf{Tier E --- Extension},
    fonttitle=\bfseries, arc=1.5mm, left=3mm, right=3mm, top=2mm, bottom=2mm, breakable]
\textbf{Part 1 --- No graph allowed.} Find each domain and vertical asymptote from the equation alone.
\begin{center}
\renewcommand{\arraystretch}{1.6}
\begin{tabularx}{\linewidth}{@{}l X X X@{}}
\hline
\textbf{Function} & \textbf{Condition on the argument} & \textbf{Domain (interval notation)} & \textbf{Vertical asymptote} \\
\hline
$\log_2(x-5)$ & \ans{$x-5>0$} & \ans{$(5,\infty)$} & \ans{$x=5$} \\
$\log_3(x+4)$ & \ans{$x+4>0$} & \ans{$(-4,\infty)$} & \ans{$x=-4$} \\
$\log(2x-6)$ & \ans{$2x-6>0$} & \ans{$(3,\infty)$} & \ans{$x=3$} \\
$\log_5(9-x)$ & \ans{$9-x>0$} & \ans{$(-\infty,9)$} & \ans{$x=9$} \\
\hline
\end{tabularx}
\end{center}
Exactly one of these four graphs runs to the \textbf{left}. Which one, and what in the equation tells
you so \emph{before} you graph anything? \ansline{$\log_5(9-x)$. Solving $9-x>0$ requires dividing by
a negative, which flips the inequality to $x<9$ --- so the legal inputs sit to the \emph{left} of the
asymptote. The tell is the $-x$: the variable is being subtracted, so the argument shrinks as $x$
grows.}

\vspace{5pt}
\textbf{Part 2 --- A logarithmic model, and the family it came from.} A bacteria colony doubles every
hour. If the colony is now $n$ times its original size, the number of hours that have passed is
\[
  t(n) = \log_2 n \ \text{ hours.}
\]
\begin{center}
\renewcommand{\arraystretch}{1.25}
\begin{tabular}{|c|c|c|c|c|c|}
\hline
$n$ (times original size) & $1$ & $2$ & $4$ & $16$ & $256$ \\ \hline
$t(n)$ (hours) & $0$ & $1$ & $2$ & $4$ & $8$ \\ \hline
\end{tabular}
\end{center}
\begin{enumerate}[label=(\alph*), leftmargin=1.9em, itemsep=6pt]
  \item What is the domain of $t$ \emph{algebraically} (from the argument)? \ans{$(0,\infty)$} \\
        Does the context agree, or does it restrict things further? \ansline{The context agrees and
        then narrows it: a colony that has been growing since time zero is at least its original
        size, so realistically $n\ge 1$, giving $t\ge 0$. Sizes between $0$ and $1$ are legal for the
        function but would describe a colony \emph{shrinking}, which this model does not do.}
  \item The graph of $t$ has an $x$-intercept at $(1,0)$. What does that point \emph{mean} for the
        colony? \ansline{$n=1$ means the colony is $1$ times its original size --- it has not grown
        at all --- and $t=0$ says no time has passed. The intercept is the starting instant.}
  \item From the table: going from $n=16$ to $n=256$ multiplies the colony size by \ans{16},
        but adds only \ans{4} hours. Use the shape of a logarithmic graph to explain why the
        time grows so much more slowly than the size. \ansline{A logarithmic graph flattens as it
        moves right: equal \emph{additions} to the output require repeated \emph{multiplications} of
        the input. Each extra hour only doubles the colony, so to add $4$ hours you must double four
        times --- a factor of $16$. Multiplying the input by a huge number moves the output by a
        small amount.}
  \item Rewrite $t=\log_2 n$ as an exponential statement solved for $n$. \quad $n=$ \ans{$2^{t}$} \\
        Which function family is that? \ans{exponential} \quad What does this tell you about the
        relationship between $t(n)$ and the function you just wrote? \ansline{They are
        \textbf{inverses}. One takes a size and returns the hours; the other takes the hours and
        returns the size. Their tables are the same table with the rows traded, and their graphs are
        reflections over $y=x$.}
  \item The model has a vertical asymptote at $n=0$. In terms of the bacteria, explain what the
        asymptote is saying --- and why the colony size can get as close to $0$ as you like without
        the model ever accepting $n=0$ itself. \ansline{Running the doubling backwards halves the
        colony each hour: $\tfrac12$, $\tfrac14$, $\tfrac18$, \dots\ The size gets arbitrarily close
        to zero but never reaches it, so the ``hours ago'' answer runs to $-\infty$ without ever
        stopping.\\[3pt]
        A size of exactly $0$ would need a power of $2$ equal to $0$, and there is none --- so $n=0$
        is outside the domain. That is precisely what a vertical asymptote records: the graph
        approaches the line $n=0$ forever and never arrives.}
\end{enumerate}
\end{tcolorbox}

\begin{teachernote}
\textbf{Tier placement.} Send students who wrote square brackets on the notes domains to Tier~R --- the
whole tier is built to force the asymptote-vs-endpoint distinction (items~2 and~4). Students who
answered the Section~3 ``which four rows are identical'' question well can start at Tier~E.

\textbf{Tier~A is the compare-and-contrast standard (A2.F.2b/A2.F.1e).} The table is designed so
five rows differ and two agree; the payoff is Justify~1, where students discover that the
$y$-intercept is not a property of ``being a log'' but of \emph{where the asymptote sits}. Expect
some groups to claim a log graph can \emph{never} have a $y$-intercept --- $p$ is the counterexample,
and it is worth pausing the room for.

\textbf{Tier~E Part~2(d)} is the first time students convert a logarithmic statement to exponential
form. Accept an informal answer read off the definition (``$2$ to the $t$ gives $n$''); the formal
two-way conversion is Lesson~8.1's opening move. Part~(e) is the strongest item on the page --- it
asks for an asymptote's \emph{meaning} in context, not its equation.
\end{teachernote}

\end{document}
